\documentclass[UTF8, 12pt, fontset=fandol, oneside]{ctexart}
\usepackage{researchnotes}

\title{强化学习：导论}
\author{}
\date{}

\begin{document}

\maketitle

\section*{第 6 章\quad 时序差分学习}

如果必须指出强化学习中一个既核心又新颖的思想，那毫无疑问就是\textbf{时序差分}（temporal-difference，TD）学习。
TD 学习是蒙特卡洛思想和动态规划（DP）思想的结合。
和蒙特卡洛方法一样，TD 方法也可以直接从原始经验中学习，而不需要环境动态的模型。
和 DP 一样，TD 方法会根据其他已学到的估计来更新估计，而不必等待最终结果（也就是说，它们会自举）。
TD、DP 和蒙特卡洛方法之间的关系，是强化学习理论中反复出现的主题；本章只是我们探索它的开始。
在结束之前，我们会看到这些思想和方法怎样彼此融合，也能怎样以许多方式组合起来。
特别地，在第 7 章中，我们将介绍 $n$ 步算法，它在 TD 到蒙特卡洛方法之间架起桥梁；在第 12 章中，我们将介绍 TD($\lambda$) 算法，它能无缝地统一这些算法。

和往常一样，我们先聚焦于策略评价或\textbf{预测}问题，也就是估计给定策略 $\pi$ 的价值函数 $v_\pi$ 的问题。
对于\textbf{控制}问题（寻找一个最优策略），DP、TD 和蒙特卡洛方法都使用某种形式的广义策略迭代（GPI）。
这些方法之间的差异，主要体现在它们处理预测问题的方式上。

\subsection*{6.1\quad TD 预测}

TD 方法和蒙特卡洛方法都使用经验来求解预测问题。
给定一些遵循策略 $\pi$ 的经验，两种方法都会更新对非终止状态 $S_t$ 中出现状态 $s$ 时 $v_\pi$ 的估计 $V$。
粗略地说，蒙特卡洛方法会等到访问之后的回报已知，再把该回报作为 $V(S_t)$ 的目标。
一个适用于非平稳环境的简单每次访问蒙特卡洛方法是
\begin{equation}
V(S_t)\leftarrow V(S_t)+\alpha\left[G_t-V(S_t)\right].
\end{equation}
其中 $G_t$ 是时刻 $t$ 之后的实际回报，$\alpha$ 是一个常数步长参数（参见式（2.4））。
我们把这个方法称为\textbf{常数-$\alpha$ MC}。
蒙特卡洛方法必须等到一幕结束，才能确定对 $V(S_t)$ 的增量（只有到那时 $G_t$ 才已知），而 TD 方法只需要等到下一个时间步。
在时刻 $t+1$，它们立刻形成一个目标，并用观测到的奖励 $R_{t+1}$ 和估计 $V(S_{t+1})$ 做出有用的更新。
最简单的 TD 方法做出的更新为
\begin{equation}
V(S_t)\leftarrow V(S_t)+\alpha\left[R_{t+1}+\gamma V(S_{t+1})-V(S_t)\right],
\end{equation}
这个更新在转移到 $S_{t+1}$ 并接收到 $R_{t+1}$ 之后立即执行。
事实上，蒙特卡洛更新的目标是 $G_t$，而 TD 更新的目标是 $R_{t+1}+\gamma V(S_{t+1})$。
这种 TD 方法称为 TD(0)，或\textbf{一步 TD}，因为它是第 7 章中 TD($\lambda$) 和 $n$ 步 TD 方法的一个特例。
下面的框给出了 TD(0) 的完整过程形式。

\begin{center}
\fbox{\begin{minipage}{0.88\textwidth}
\noindent\textbf{表格型 TD(0)，用于估计 $v_\pi$}

\vspace{0.4em}
\noindent 输入：要评价的策略 $\pi$。

\noindent 算法参数：步长 $\alpha\in(0,1]$。

\noindent 初始化 $V(s)$，对所有 $s\in\mathcal{S}^+$ 任意；但 $V(terminal)=0$。

\noindent 对每一幕循环：

\quad 初始化 $S$。

\quad 对一幕中的每一步循环：

\qquad 按照 $\pi$ 在 $S$ 中采取动作 $A$。

\qquad 采取动作 $A$，观察 $R,S'$。

\qquad $V(S)\leftarrow V(S)+\alpha\left[R+\gamma V(S')-V(S)\right]$。

\qquad $S\leftarrow S'$。

\quad 直到 $S$ 是终止状态。
\end{minipage}}
\end{center}

因为 TD(0) 的更新部分地基于一个已有估计，所以我们说它是一种\textbf{自举}方法，和 DP 一样。
由第 3 章可知
\begin{equation}
\begin{aligned}
v_\pi(s)&=\mathbb{E}_\pi[G_t\mid S_t=s]\\
&=\mathbb{E}_\pi[R_{t+1}+\gamma G_{t+1}\mid S_t=s]\qquad \text{（由（3.9））}\\
&=\mathbb{E}_\pi[R_{t+1}+\gamma v_\pi(S_{t+1})\mid S_t=s]\qquad \text{（6.4）}.
\end{aligned}
\end{equation}
粗略地说，蒙特卡洛方法用式（6.3）的一个估计作为目标，而 DP 方法用式（6.4）的一个估计作为目标。
蒙特卡洛目标是一个估计，因为式（6.3）中的期望并不知道；用的是一个样本回报来代替真实期望回报。
DP 目标不是估计，因为期望值由环境模型完整给出；但因为 $v_\pi(S_{t+1})$ 不知道，所以 DP 使用当前估计 $V(S_{t+1})$ 来代替。
TD 目标之所以是估计，两个原因兼而有之：它对式（6.4）中的期望值进行采样，并且用当前估计 $V$ 代替真实的 $v_\pi$。
因此，TD 方法把蒙特卡洛的采样与 DP 的自举结合起来。
正如我们将看到的，只要足够谨慎而有想象力，这能让我们在很大程度上同时获得蒙特卡洛方法和 DP 方法的优势。

右侧所示的是表格型 TD(0) 的备份图。
被更新的状态的价值估计位于备份图的顶端；更新所依据的是从该状态到紧随其后的状态的一次样本转移。
我们把 TD 和蒙特卡洛更新称为\textbf{样本更新}，因为它们向前看的是一个样本后继状态（或状态--动作对），使用后继状态的价值以及沿途的奖励来计算备份值，然后相应地更新原状态（或状态--动作对）的价值。
\textbf{样本更新}不同于 DP 方法中的\textbf{期望更新}，后者基于所有可能后继状态的完整分布。

\begin{figure}[htbp]
\centering
\includegraphics[height=0.18\textheight]{tex/强化学习/images/math-7/math-7-td0-backup.png}
\end{figure}

最后，注意 TD(0) 更新中方括号内的量是一种误差，它衡量 $S_t$ 的估计价值与更好估计 $R_{t+1}+\gamma V(S_{t+1})$ 之间的差。
这个量称为 \textbf{TD 误差}，在强化学习中会以多种形式出现：
\begin{equation}
\delta_t\doteq R_{t+1}+\gamma V(S_{t+1})-V(S_t).
\end{equation}
注意，时刻 $t$ 的 TD 误差，是\textbf{在该时刻形成的估计}中的误差。
因为 TD 误差依赖下一个状态和下一个奖励，所以直到一个时间步之后才真正可用。
也就是说，$\delta_t$ 是 $t+1$ 时刻可用的 $V(S_t)$ 中的误差。
还要注意，如果在一幕中数组 $V$ 不发生变化（如蒙特卡洛方法中那样），那么蒙特卡洛误差可以写成 TD 误差之和：
\begin{equation}
\begin{aligned}
G_t-V(S_t)
&=R_{t+1}+\gamma G_{t+1}-V(S_t)+\gamma V(S_{t+1})-\gamma V(S_{t+1})
\qquad \text{（由（3.9））}\\
&=\delta_t+\gamma\left(G_{t+1}-V(S_{t+1})\right)\\
&=\delta_t+\gamma\delta_{t+1}+\gamma^2\left(G_{t+2}-V(S_{t+2})\right)\\
&=\delta_t+\gamma\delta_{t+1}+\gamma^2\delta_{t+2}+\cdots+\gamma^{T-t-1}\delta_{T-1}
  +\gamma^{T-t}\left(G_T-V(S_T)\right)\\
&=\delta_t+\gamma\delta_{t+1}+\gamma^2\delta_{t+2}+\cdots+\gamma^{T-t-1}\delta_{T-1}
  +\gamma^{T-t}(0-0)\\
&=\sum_{k=t}^{T-1}\gamma^{k-t}\delta_k .
\end{aligned}
\end{equation}
如果在一幕中像 TD(0) 那样更新 $V$，这个恒等式就不精确；但如果步长很小，它仍可能近似成立。
这个恒等式的推广在资格迹和时序差分学习的理论与算法中起着重要作用。

\noindent\textbf{练习 6.1}
如果 $V$ 在一幕中发生变化，那么式（6.6）只近似成立；精确的差别会是什么？
设 $V_t$ 表示在 TD 误差（6.5）和 TD 更新（6.2）中于时刻 $t$ 使用的状态价值数组。
在上面的推导中把额外项重新推导出来，以确定必须加到 TD 误差之和上的量，使其等于蒙特卡洛误差。

\noindent\textbf{例 6.1：开车回家}
每天你下班开车回家时，都会试图预测还要多久才能到家。
离开办公室时，你会记下星期几、当天天气，以及任何可能相关的情况。
假设今天是星期五，你正好在 6 点离开，并估计回家需要 30 分钟。
你在 6:05 到达车旁，注意到开始下雨。
雨天交通通常更慢，所以你把从现在开始还需的时间重新估计为 35 分钟，也就是总共 40 分钟。
15 分钟后，你已经跑完一段高速公路，而你的预测没有变。
又过 10 分钟，你离开高速公路进入一条支路，这时你把总旅行时间缩短为 35 分钟。
遗憾的是，就在这时你被一辆慢车堵住，前面道路太窄无法超车。
你最终在 6:40 转到自己住处所在的支路。
3 分钟后，你到家了。
下面给出了各状态、时间和预测的序列：

\begin{center}
\begin{tabular}{lccc}
\toprule
状态 & 已用时间（分钟） & 预测剩余时间 & 预测总时间 \\
\midrule
星期五 6 点离开办公室 & 0 & 30 & 30 \\
到达车旁，开始下雨 & 5 & 35 & 40 \\
驶离高速公路 & 20 & 15 & 35 \\
第 2 条路，堵在卡车后 & 30 & 10 & 40 \\
进入住处所在支路 & 40 & 3 & 43 \\
到家 & 43 & 0 & 43 \\
\bottomrule
\end{tabular}
\end{center}

这个例子中的奖励，是每段旅程上经过的时间。\footnote{如果我们关心的不是客观的最小旅行时间，而是控制问题，那么当然会把奖励设为经过时间的负值。不过这里我们只关心预测（策略评价），因此可以把符号反过来，用正数来简化表达。}
我们不折扣（$\gamma=1$），于是每个状态的回报就是从该状态到目的地的实际剩余时间。
最后一列给出了每个遇到的状态的当前估计价值。

观察蒙特卡洛方法的运行方式，一个简单办法是把这一序列上的预测总时间（最后一列）画出来，如图 6.1（左）所示。
红色箭头表示常数-$\alpha$ MC 方法在 $\alpha=1$ 时建议的预测变化。
这些变化正是每个状态的估计价值（预测剩余时间）与实际回报（实际剩余时间）之间的误差。
例如，在驶离高速公路时，你以为还只需要 15 分钟就能到家，但事实上需要 23 分钟。
式（6.1）在这种情况下给出的增量，是把驶离高速公路之后的时间估计增加 8 分钟。
误差 $G_t-V(S_t)$ 此时是 8 分钟。
假设步长参数 $\alpha$ 为 $1/2$，那么同一个预测时间只会被改变 4 分钟。
请注意，蒙特卡洛方法要等到最终结果已知后才能学习，也就是必须等到你到家。
这在这里大概是一个过大的变化；卡车很可能只是一次偶然拥堵。
无论如何，变化只能在离线方式中进行，也就是在你到家之后；只有到那时你才知道任何实际回报。

是否必须等到最终结果已知后才能开始学习？
假设另一天，你离开办公室时同样估计需要 30 分钟才能开回家，但之后遇到严重堵车。
离开办公室 20 分钟后，你仍在高速路上走走停停。
现在你估计还要 25 分钟才能到家，总共 50 分钟。
等待在交通中时，你已经知道最初 30 分钟的估计过于乐观。
是否必须等到回家之后，才能增加对初始状态的估计？
按照蒙特卡洛方法，你必须等，因为你还不知道回报。

按照 TD 方法则不然，你会立即学习，把你的初始估计从 30 分钟朝 50 分钟移动。
事实上，每次估计都将向紧随其后的估计移动。
回到图 6.1，右侧显示了 TD 规则（6.2）在 $\alpha=1$ 时建议的预测变化。
这些变化中每一个都正比于随时间变化的预测，也就是预测中的\textbf{时序差分}。

\begin{figure}[htbp]
\centering
\includegraphics[width=0.82\textwidth]{tex/强化学习/images/math-7/math-7-fig-6-1.png}
\caption{在开车回家例子中，蒙特卡洛方法（左）和 TD 方法（右）建议的变化。}
\end{figure}

除了可以在等待交通时做事之外，TD 学习还有其他计算上的好处：基于当前预测而非等到终止、实际回报已知后再学习，是有利的。
下一节将对此做一个简要讨论。

\noindent\textbf{练习 6.2}
这个练习帮助你发展关于 TD 方法为何常常比蒙特卡洛方法更高效的直觉。
考虑开车回家例子，以及 TD 方法和蒙特卡洛方法如何处理它。
你能否想象一种场景，使得 TD 更新平均而言会优于蒙特卡洛更新？
请给出一个样例场景，也就是一个过去经验和当前状态的描述，使你预期 TD 更新会更好。
提示：假设你有大量从工作地点开车回家的经验。
于是你搬到了新家，开始学习从工作地点开车到新家的预测。
现在你正在学的预测是从工作地点到家的时间。
你能看出为什么 TD 更新很可能会优于蒙特卡洛更新吗？
起初，在这种情形下情况是怎样的？
原来的场景中是否也可能发生同类事情？

\subsection*{6.2\quad TD 预测方法的优点}

TD 方法部分地基于其他估计来更新自己的估计。
它们从一个猜测中学习，也就是说，它们会自举。
这样做是好事吗？
TD 方法相对于蒙特卡洛方法和 DP 方法有什么优势？
发展并回答这些问题，将占据本书余下部分的许多内容。
在本节中，我们简要预览其中一些答案。

显然，TD 方法相对于 DP 方法的一个优势是，它们不需要环境模型，不需要奖励和下一状态概率分布。

TD 方法相对于蒙特卡洛方法的下一个最明显优势，是它们天然以在线、完全增量的方式实现。
使用蒙特卡洛方法时，必须等到一幕结束，因为只有那时回报才已知；而使用 TD 方法时，只需等待一个时间步。
令人惊讶的是，这常常是一个关键考虑。
有些应用的幕很长，以至于把所有学习都推迟到一幕结束太慢。
另一些应用是持续任务，根本没有幕。
最后，正如我们在上一章所指出的，某些蒙特卡洛方法必须忽略或折扣那些采取了实验性动作的幕，这会极大地减慢学习。
TD 方法受这些问题影响小得多，因为无论后续采取了什么动作，它们都从每一次转移中学习。

但是，TD 方法是否可靠？
当然，从下一个猜测中学习而不等待实际结果，是很方便的；但它是否仍能保证收敛到正确答案？
幸运的是，答案是肯定的。
对于任意固定策略 $\pi$，如果步长参数是常数且足够小，那么 TD(0) 已被证明会以均值形式收敛到 $v_\pi$；如果步长参数按通常的随机逼近条件（2.7）递减，则以概率 1 收敛。
大多数收敛证明只适用于上面给出的算法（6.2）的表格型情形，但也有一些证明适用于一般线性函数近似情形。
这些结果将在第 9 章更一般地讨论。

如果 TD 方法和蒙特卡洛方法都渐近收敛到正确预测，那么自然的下一个问题是：“谁先到达那里？”
换句话说，哪种方法学得更快？
哪种方法更有效地利用有限数据？
目前，这在某种意义上仍是一个开放问题，因为还没有人能够从数学上证明其中一种方法比另一种收敛得更快。
事实上，甚至怎样最恰当地从形式上表述这个问题也不完全清楚。
不过在实践中，如例 6.2 所示，在随机任务上，TD 方法通常被发现比常数-$\alpha$ MC 方法收敛更快。

\noindent\textbf{例 6.2：随机游走}
在这个例子中，我们经验性地比较 TD(0) 和常数-$\alpha$ MC 在下面这个马尔可夫奖励过程中的预测能力：

\begin{figure}[htbp]
\centering
\includegraphics[width=0.72\textwidth]{tex/强化学习/images/math-7/math-7-random-walk-chain.png}
\end{figure}

马尔可夫奖励过程（MRP）是不含动作的马尔可夫决策过程。
我们常会在关注预测问题时使用 MRP；在这种情形下，无需区分环境动态和由智能体造成的动态。
在这个 MRP 中，所有幕都从中心状态 $C$ 开始，然后每一步以相等概率向左或向右移动一个状态。
一幕在最左端或最右端终止。
当一幕在右端终止时，得到 $+1$ 的奖励；其他所有奖励为零。
例如，一个典型幕可能由如下状态--奖励序列组成：$C,0,B,0,C,0,D,0,E,1$。
因为这个任务不折扣，所以每个状态的真实价值，就是从该状态开始终止在右端的概率。
因此，中心状态 $C$ 的真实价值是 $v_\pi(C)=0.5$。
从 $A$ 到 $E$，所有状态的真实价值分别为 $\frac16,\frac26,\frac36,\frac46,\frac56$。

\begin{figure}[htbp]
\centering
\includegraphics[width=0.86\textwidth]{tex/强化学习/images/math-7/math-7-random-walk-plots.png}
\end{figure}

上面的左图显示了 TD(0) 的一次运行在若干幕数之后学到的价值。
100 幕之后的估计大致已经接近真实价值；随后在常数步长参数的作用下（本例中 $\alpha=0.1$），这些价值会随着最近若干幕的结果而无限波动。
右图显示了两种方法在不同 $\alpha$ 值下的学习曲线。
所示性能指标是所学价值函数与真实价值函数之间的均方根（RMS）误差，先在五个状态上平均，再在 100 次运行上平均。
在所有情形下，近似价值函数都被初始化为中间值 $V(s)=0.5$，对所有 $s$ 相同。
在这个任务上，TD 方法始终优于 MC 方法。

\noindent\textbf{练习 6.3}
从随机游走例子左图所示结果看，第一幕后只改变了 $V(A)$ 的估计。
这说明第一幕发生了什么？
为什么只有这个状态的估计被改变？
究竟被改变了多少？

\noindent\textbf{练习 6.4}
随机游走例子右图中的具体结果依赖于步长参数 $\alpha$ 的取值。
你是否认为，如果使用更宽范围的 $\alpha$ 值，对哪种算法更好的结论会受影响？
是否存在另一个固定的 $\alpha$ 值，使得某个算法的表现会显著优于图中所示？
为什么或为什么不？

\noindent\textbf{练习 6.5}
在随机游走例子的右图中，TD 方法的 RMS 误差似乎先下降后上升，特别是在较大的 $\alpha$ 下。
是什么可能造成这种现象？
你认为这总会发生，还是可能取决于近似价值函数怎样初始化？

\noindent\textbf{练习 6.6}
在例 6.2 中，我们说随机游走例子的真实价值从 $A$ 到 $E$ 分别是 $\frac16,\frac26,\frac36,\frac46,\frac56$。
请描述至少两种不同方法来计算这些值。
你猜我们实际用的是哪一种？
为什么？

\subsection*{6.3\quad TD(0) 的最优性}

假设只有有限量的经验可用，例如 10 幕或 100 个时间步。
在这种情形下，一种常见的增量学习方法，是把经验重复呈现给学习方法，直到该方法收敛到一个答案。
给定一个近似价值函数 $V$，式（6.1）或式（6.2）中指定的增量会在每次访问非终止状态时计算出来，但价值函数只会在所有增量之和的基础上整体改变。
然后，用新的价值函数再次处理全部可用经验，以产生新的整体增量，如此继续，直到价值函数收敛。
我们把这种\textbf{批量更新}称为在处理完每一整批训练数据之后才进行更新。

在批量更新下，只要 $\alpha$ 选得足够小，TD(0) 会确定性地收敛到一个与步长参数 $\alpha$ 无关的单一答案。
常数-$\alpha$ MC 方法也会在相同条件下确定性收敛，但会收敛到另一个答案。
理解这两个答案，将帮助我们理解两种方法之间的差异。
在通常更新下，这些方法不会一路移动到各自的批量答案，而只是朝这些方向迈步。
在试图理解两种答案之前，先看几个例子。

\noindent\textbf{例 6.3：批量更新下的随机游走}
把 TD(0) 和常数-$\alpha$ MC 的批量更新版本应用于随机游走预测例子（例 6.2），方法如下。
每一新幕结束后，至此为止的所有幕都被视作一批。
它们被反复呈现给算法，也就是 TD(0) 或常数-$\alpha$ MC，并使用足够小的 $\alpha$，直到价值函数收敛。
然后将得到的价值函数与 $v_\pi$ 比较，并把五个状态上的平均均方误差（以及整个实验 100 次独立重复的平均）绘出，得到图 6.2 所示学习曲线。
注意，批量 TD 方法始终优于批量蒙特卡洛方法。

在批量训练下，常数-$\alpha$ MC 收敛到数值 $V(s)$，它们是访问每个状态 $s$ 后实际回报的样本平均。
这些估计是最优的，因为它们在训练集上使实际回报的均方误差最小。
从这个意义上看，批量 TD 方法能按照右图所示的均方根误差指标表现得更好，是令人惊讶的。
TD 方法怎么会比这种最优方法更好？
答案是，蒙特卡洛方法只是在一种有限意义上是最优的，而且这种意义与预测回报并不最相关。

\begin{figure}[htbp]
\centering
\includegraphics[width=0.52\textwidth]{tex/强化学习/images/math-7/math-7-fig-6-2.png}
\caption{批量训练下，TD(0) 与常数-$\alpha$ MC 在随机游走任务上的表现。}
\end{figure}

\noindent\textbf{例 6.4：你是预测器}
请把自己放在一个未知马尔可夫奖励过程的回报预测器位置上。
假设你观察到下面八个幕：
\[
\begin{array}{cc}
A,0,B,0 & B,1\\
B,1 & B,1\\
B,1 & B,1\\
B,1 & B,0
\end{array}
\]
这意味着第一幕从状态 $A$ 开始，带着奖励 $0$ 转移到 $B$，然后从 $B$ 带着奖励 $0$ 终止。
另外七个幕更短：从 $B$ 开始并立即终止。
给定这一批数据，你会说 $V(A)$ 和 $V(B)$ 的最优预测值分别是多少？
大家大概都会同意，$V(B)$ 的最优值是 $\frac34$，因为 $B$ 过程终止的八次中有六次回报为 1，另外两次回报为 0。

但是，给定这些数据，$V(A)$ 的最优值是多少？
这里有两个合理答案。
一个答案是观察到：从 $A$ 开始的过程有 100\% 的次数立即转移到 $B$（奖励为 0）；而我们已经判定 $B$ 的价值为 $\frac34$；因此 $A$ 的价值也必须为 $\frac34$。
看待这个答案的一种方式是，它基于先对马尔可夫过程建模，在这个例子中模型如下右图所示，然后在给定模型的情况下计算正确估计，这个模型给出的 $V(A)=\frac34$。
这也是批量 TD(0) 给出的答案。

\begin{figure}[htbp]
\centering
\includegraphics[width=0.42\textwidth]{tex/强化学习/images/math-7/math-7-predictor-mrp.png}
\end{figure}

另一个合理答案是简单观察到：我们见过一次 $A$，随后得到的回报是 0；因此估计 $V(A)$ 为 0。
这正是批量蒙特卡洛方法给出的答案。
请注意，这也是在训练数据上给出最小均方误差的答案。
事实上，它在数据上给出零误差。
但我们仍然预期第一个答案更好。
如果这个过程是马尔可夫的，我们就预期第一个答案在\textbf{未来}数据上会产生更低误差，即使蒙特卡洛答案在已有数据上更好。

例 6.4 说明了批量 TD(0) 和批量蒙特卡洛方法得到的估计之间的一般差异。
批量蒙特卡洛方法总是找到使训练集上均方误差最小的估计；而批量 TD(0) 总是找到对马尔可夫过程最大似然模型而言恰好正确的估计。
一般来说，最大似然估计是参数值，其生成数据的概率最大。
在这个例子中，最大似然估计显然是以显而易见的方式从观测幕形成的马尔可夫过程模型：从 $i$ 到 $j$ 的估计转移概率，是从 $i$ 到 $j$ 的观测转移次数占从 $i$ 出发的全部观测转移次数的比例；相关的期望奖励，是这些转移上观测奖励的平均值。
给定这个模型，我们就可以计算价值函数的估计，这个估计在模型完全正确时会是完全正确的。
这称为\textbf{确定性等价估计}，因为它等价于假设底层过程的估计近似是确定的，而不是近似的。
一般而言，批量 TD(0) 会收敛到确定性等价估计。

这有助于解释为什么 TD 方法会比蒙特卡洛方法收敛得更快。
在批量形式下，TD(0) 比蒙特卡洛方法更快，因为它计算真实的确定性等价估计。
这解释了图 6.2 中批量随机游走任务上 TD(0) 的优势。
与确定性等价估计的关系，也可能部分解释非批量 TD(0) 的速度优势（例如例 6.2 第 125 页右图）。
尽管非批量方法既不能达到确定性等价估计，也不能达到最小均方误差估计，但可以把它们大致理解为沿着这些方向移动。
非批量 TD(0) 可能比常数-$\alpha$ MC 更快，因为它朝着一个更好的估计移动，即使它没有一路到达那里。
目前，对于在线 TD 和蒙特卡洛方法的相对效率，还没有更明确的说法。

最后，值得注意的是，虽然确定性等价估计在某种意义上是一个最优解，但几乎从来无法直接计算它。
如果 $n=|\mathcal{S}|$ 是状态数，那么仅形成过程的最大似然估计就可能需要 $n^2$ 阶的内存，而以常规方式计算相应价值函数则需要 $n^3$ 阶的计算步骤。
用这些术语来看，TD 方法能够使用不超过 $n$ 阶的内存，并在训练集上重复计算不超过 $n$ 阶的时间内近似同样的解，确实令人印象深刻。
在状态空间很大的任务上，TD 方法可能是近似确定性等价解的唯一可行方式。

\noindent\textbf{练习 6.7}
设计一种 TD(0) 更新的离策略版本，使其可用于
任意目标策略 $\pi$ 和覆盖它的行为策略 $b$，并在每一步使用重要性采样比 $\rho_{t:t}$（5.3）。

\subsection*{6.4\quad Sarsa：同策略 TD 控制}

现在我们转向使用 TD 预测方法来处理控制问题。
和往常一样，我们遵循广义策略迭代（GPI）的模式，只是这一次把 TD 方法用于评价或预测部分。
和蒙特卡洛方法一样，我们面临探索与利用之间的权衡，并且同样把方法分成两大类：同策略和离策略。
本节介绍一种同策略 TD 控制方法。

第一步是学习一个动作价值函数，而不是状态价值函数。
特别地，对于同策略方法，我们必须估计当前行为策略 $\pi$ 以及所有状态 $s$ 和动作 $a$ 的 $q_\pi(s,a)$。
这可以用本章前面对学习 $v_\pi$ 所描述的同一个 TD 方法完成。
回忆一幕由状态和状态--动作对交替构成的序列组成：
\[
\cdots,S_t,A_t,R_{t+1},S_{t+1},A_{t+1},R_{t+2},S_{t+2},A_{t+2},R_{t+3},S_{t+3},A_{t+3},\cdots
\]
上一节中，我们考察从状态到状态的转移，并学习状态的价值。
现在，我们考察从状态--动作对到状态--动作对的转移，并学习状态--动作对的价值。
形式上，这两种情形是相同的：二者都是带奖励过程的马尔可夫链。
状态价值在 TD(0) 下收敛的定理，也适用于对应的动作价值算法：
\begin{equation}
Q(S_t,A_t)\leftarrow Q(S_t,A_t)+\alpha\left[R_{t+1}+\gamma Q(S_{t+1},A_{t+1})-Q(S_t,A_t)\right].
\end{equation}
这个更新在每次从非终止状态 $S_t$ 转移之后完成。
如果 $S_{t+1}$ 是终止状态，那么 $Q(S_{t+1},A_{t+1})$ 被定义为零。
这条规则使用五元组事件 $(S_t,A_t,R_{t+1},S_{t+1},A_{t+1})$ 中的每一个元素。
这个五元组从一个状态--动作对转移到下一个状态--动作对，于是给算法起名为 Sarsa。
其备份图如下所示。

\begin{figure}[htbp]
\centering
\includegraphics[height=0.18\textheight]{tex/强化学习/images/math-7/math-7-sarsa-backup.png}
\end{figure}

基于 Sarsa 预测方法来设计一个同策略控制算法是直接的。
和所有同策略方法一样，我们持续估计行为策略 $\pi$ 的 $q_\pi$，同时又把 $\pi$ 朝关于 $q_\pi$ 贪心的方向改变。
Sarsa 控制算法的一般形式见下一页的框。

Sarsa 算法的收敛性质，依赖于策略对 $Q$ 的依赖性质。
例如，可以使用 $\varepsilon$-贪心或 $\varepsilon$-soft 策略。
如果所有状态--动作对都被无限次访问，并且策略在极限中收敛到贪心策略，那么 Sarsa 以概率 1 收敛到最优策略和最优动作价值函数；例如，可以令 $\varepsilon$-贪心策略中的 $\varepsilon=1/t$。

\noindent\textbf{练习 6.8}
证明动作价值形式的 TD 误差
\[
\delta_t=R_{t+1}+\gamma Q(S_{t+1},A_{t+1})-Q(S_t,A_t)
\]
同样满足式（6.6）的动作价值版本，仍假设价值在各步之间不发生改变。

\begin{center}
\fbox{\begin{minipage}{0.88\textwidth}
\noindent\textbf{Sarsa（同策略 TD 控制），用于估计 $Q\approx q_*$}

\vspace{0.4em}
\noindent 算法参数：步长 $\alpha\in(0,1]$，小的 $\varepsilon>0$。

\noindent 初始化 $Q(s,a)$，对所有 $s\in\mathcal{S}^+$、$a\in\mathcal{A}(s)$ 任意；但 $Q(terminal,\cdot)=0$。

\noindent 对每一幕循环：

\quad 初始化 $S$。

\quad 用由 $Q$ 导出的策略（例如 $\varepsilon$-贪心）从 $S$ 选择 $A$。

\quad 对一幕中的每一步循环：

\qquad 采取动作 $A$，观察 $R,S'$。

\qquad 用由 $Q$ 导出的策略（例如 $\varepsilon$-贪心）从 $S'$ 选择 $A'$。

\qquad $Q(S,A)\leftarrow Q(S,A)+\alpha\left[R+\gamma Q(S',A')-Q(S,A)\right]$。

\qquad $S\leftarrow S'$；$A\leftarrow A'$。

\quad 直到 $S$ 是终止状态。
\end{minipage}}
\end{center}

\noindent\textbf{例 6.5：有风的网格世界}
下图插入处显示的是一个标准网格世界，包含起点和目标状态，但有一个不同之处：网格中间有一股向上的侧风。
动作是标准的四个动作：上、下、右、左；但在中间区域，结果的下一状态会被“风”向上推动，推动强度随列而变。
风的强度在每列下面给出，单位为向上移动的格数。
例如，如果你在目标右侧一格，那么动作 left 会把你带到目标正上方一格。
这是一个未折扣的分幕任务，奖励恒为 $-1$，直到到达目标状态。

\begin{figure}[htbp]
\centering
\includegraphics[width=0.58\textwidth]{tex/强化学习/images/math-7/math-7-windy-gridworld.png}
\end{figure}

右侧曲线显示了把 $\varepsilon$-贪心 Sarsa 应用于这个任务的结果，其中 $\varepsilon=0.1$、$\alpha=0.5$，初始价值 $Q(s,a)=0$，对所有 $s,a$ 都如此。
图中斜率增加表明，随着时间推移，目标达成得越来越快。
到 8000 个时间步时，贪心策略早已达到最优（图中插入处给出了一条轨迹）；继续进行的 $\varepsilon$-贪心探索使得平均幕长度保持在约 17 步，比最小值 15 只多约两步。
请注意，不能轻易在这里使用蒙特卡洛方法，因为并不保证所有策略都会终止。
如果一旦发现某个策略会使智能体一直停留在同一状态，那么下一幕将永远不会结束。
像 Sarsa 这样的在线学习方法没有这个问题，因为它们会在一幕中很快学到这类策略很差，并改换到其他策略。

\noindent\textbf{练习 6.9：有国王走法的有风网格世界（编程）}
重新求解有风网格世界，假设可行动作中包括对角移动，而不只是通常的四个动作。
额外动作能带来多少改进？
如果还包括一个不移动动作，你还能做得更好吗？
在被风推离目标时，不移动动作怎样帮助你保持在目标处？

\noindent\textbf{练习 6.10：随机风（编程）}
重新求解带国王走法的有风网格世界任务，但假设风的影响如果存在，则是随机的，有时比每列给出的平均值大 1，有时小 1。
也就是说，在上一练习中，你会严格按给出的值移动；但在这个练习中，有三分之一时间向上多移动一格，三分之一时间向上少移动一格。
例如，如果你在目标右侧一格并向左移动，那么有三分之一时间你会移动到目标正上方一格，三分之一时间会移动到目标上方两格，三分之一时间会移动到目标。

\subsection*{6.5\quad Q-learning：离策略 TD 控制}

强化学习早期突破之一，是发展出一种称为 \textbf{Q-learning} 的离策略 TD 控制算法（Watkins, 1989），其定义为
\begin{equation}
Q(S_t,A_t)\leftarrow Q(S_t,A_t)+\alpha\left[R_{t+1}+\gamma\max_a Q(S_{t+1},a)-Q(S_t,A_t)\right].
\end{equation}
在这种情形下，学到的动作价值函数 $Q$ 直接近似 $q_*$，也就是被遵循策略之外的最优动作价值函数。
这极大地简化了算法分析，并使早期收敛证明成为可能。
策略仍然有影响，因为它决定哪些状态--动作对被访问和更新。
然而，为了正确收敛，所需的只是所有状态--动作对都持续被更新。
正如第 5 章所观察到的，在一般情形中，任何保证找到最优行为的方法似乎也必须满足这个极小要求。
在这个假设以及通常的随机逼近条件下，$Q$ 已被证明会以概率 1 收敛到 $q_*$。
Q-learning 算法以下面的过程形式给出。

\begin{center}
\fbox{\begin{minipage}{0.88\textwidth}
\noindent\textbf{Q-learning（离策略 TD 控制），用于估计 $\pi\approx\pi_*$}

\vspace{0.4em}
\noindent 算法参数：步长 $\alpha\in(0,1]$，小的 $\varepsilon>0$。

\noindent 初始化 $Q(s,a)$，对所有 $s\in\mathcal{S}^+$、$a\in\mathcal{A}(s)$ 任意；但 $Q(terminal,\cdot)=0$。

\noindent 对每一幕循环：

\quad 初始化 $S$。

\quad 对一幕中的每一步循环：

\qquad 用由 $Q$ 导出的策略（例如 $\varepsilon$-贪心）从 $S$ 选择 $A$。

\qquad 采取动作 $A$，观察 $R,S'$。

\qquad $Q(S,A)\leftarrow Q(S,A)+\alpha\left[R+\gamma\max_a Q(S',a)-Q(S,A)\right]$。

\qquad $S\leftarrow S'$。

\quad 直到 $S$ 是终止状态。
\end{minipage}}
\end{center}

Q-learning 的备份图是什么？
规则（6.8）更新一个状态--动作对，所以备份图的顶端节点，即更新的根，必须是一个小的填充动作节点。
更新还\textbf{来自}动作节点，对下一状态中所有可能动作取最大。
因此，备份图底部节点应当全都是动作节点。
最后，请记住，我们用穿过这些“下一动作”节点的一条弧来表示对它们取最大（图 3.4 右）。
你能猜出这张图现在会是什么样吗？
如果可以，请先做出猜测，再看第 134 页图 6.4 中给出的答案。

\noindent\textbf{例 6.6：悬崖行走}
这个网格世界例子比较 Sarsa 和 Q-learning，突出同策略方法（Sarsa）与离策略方法（Q-learning）之间的差别。
考虑下图所示网格世界。
这是一个标准的未折扣分幕任务，有起点和目标状态，通常动作为上、下、右、左。
所有转移的奖励都是 $-1$，但进入标为“悬崖”的区域的奖励是 $-100$，并且会把智能体立即送回起点。

\begin{figure}[htbp]
\centering
\includegraphics[width=0.68\textwidth]{tex/强化学习/images/math-7/math-7-cliff-walking.png}
\end{figure}

右图显示了在 $\varepsilon$-贪心动作选择（$\varepsilon=0.1$）下，Sarsa 和 Q-learning 方法的表现。
经过一次初始过渡后，Q-learning 学到了最优策略的价值，也就是沿着悬崖边缘行走。
不幸的是，这会导致它偶尔因为 $\varepsilon$-贪心动作选择而跌下悬崖。
另一方面，Sarsa 会把动作选择考虑进来，并学到通过网格上半部分的较长但更安全路径。
虽然 Q-learning 实际学到的是最优策略的价值，但它的在线表现比学到迂回策略的 Sarsa 更差。
当然，如果 $\varepsilon$ 逐渐减小，那么两种方法都会渐近收敛到最优策略。

\noindent\textbf{练习 6.11}
为什么 Q-learning 被认为是一种\textbf{离策略}控制方法？

\noindent\textbf{练习 6.12}
假设动作选择是贪心的。
那么 Q-learning 是否就与 Sarsa 完全相同？
它们会做出完全相同的动作选择和权重更新吗？

\subsection*{6.6\quad Expected Sarsa}

考虑这样一种学习算法：它和 Q-learning 完全相同，只是不用对下一状态--动作对取最大，而是使用期望值，也就是考虑当前策略下每个动作的可能性。
也就是说，考虑更新规则
\begin{equation}
\begin{aligned}
Q(S_t,A_t)
&\leftarrow Q(S_t,A_t)+\alpha\left[R_{t+1}
  +\gamma\mathbb{E}_\pi[Q(S_{t+1},A_{t+1})\mid S_{t+1}]
  -Q(S_t,A_t)\right]\\
&\leftarrow Q(S_t,A_t)+\alpha\left[R_{t+1}
  +\gamma\sum_a\pi(a\mid S_{t+1})Q(S_{t+1},a)
  -Q(S_t,A_t)\right],
\end{aligned}
\end{equation}
但除此之外遵循 Q-learning 的方案。
给定下一状态 $S_{t+1}$，这个算法会确定性地朝与 Sarsa 在期望意义上相同的方向移动；因此，它称为 \textbf{Expected Sarsa}。
它的备份图在图 6.4 右侧给出。

Expected Sarsa 在计算上比 Sarsa 更复杂，但作为回报，它消除了由于随机选择 $A_{t+1}$ 而产生的方差。
给定相同数量的经验，我们可能预期它会略好于 Sarsa；事实上，它通常如此。
图 6.3 显示了悬崖行走任务上 Sarsa、Q-learning 和 Expected Sarsa 的汇总结果。
Expected Sarsa 保留了 Sarsa 相对于 Q-learning 的显著优势。
此外，在宽范围的步长参数 $\alpha$ 上，Expected Sarsa 显示出显著改进。

\begin{figure}[htbp]
\centering
\includegraphics[width=0.76\textwidth]{tex/强化学习/images/math-7/math-7-fig-6-3.png}
\caption{TD 控制方法在悬崖行走任务上的中期和渐近表现，作为 $\alpha$ 的函数。所有算法都使用 $\varepsilon$-贪心策略，其中 $\varepsilon=0.1$。渐近表现是超过 100000 幕的平均值；中期表现是前 100 幕的平均值。这些数据在中期和渐近情形下分别是超过 50000 次和 10 次运行的平均。实心圆分别标出了每种方法的最佳中期表现。改编自 van Seijen 等（2009）。}
\end{figure}

\begin{figure}[htbp]
\centering
\includegraphics[width=0.46\textwidth]{tex/强化学习/images/math-7/math-7-fig-6-4.png}
\caption{Q-learning 和 Expected Sarsa 的备份图。}
\end{figure}

这种改进在很宽范围的步长参数 $\alpha$ 上都超过了 Sarsa。
在悬崖行走中，状态转移都是确定性的，所有随机性都来自策略。
在这种情况下，Expected Sarsa 可以安全地设置 $\alpha=1$，且不会降低渐近性能；而 Sarsa 只有在较小 $\alpha$ 下才能在长期表现良好，此时短期表现很差。
在这个例子和其他例子中，都存在 Expected Sarsa 相对于 Sarsa 的一致经验优势。

在这些悬崖行走结果中，Expected Sarsa 是同策略使用的；但一般来说，它也可以使用不同于目标策略 $\pi$ 的策略来生成行为，在这种情况下，它就成为离策略算法。
例如，假设 $\pi$ 是贪心策略，而行为更具探索性；那么 Expected Sarsa 就恰好是 Q-learning。
从这个意义上说，Expected Sarsa 包含并推广了 Q-learning，同时可靠地优于 Sarsa。
除了少量额外计算成本外，Expected Sarsa 可能完全支配其他更知名的 TD 控制算法。

\subsection*{6.7\quad 最大化偏差与双重学习}

到目前为止，我们讨论过的所有控制算法，都在构造目标策略时涉及最大化。
例如，在 Q-learning 中，目标策略是给定当前动作价值下的贪心策略，而这个策略由最大化操作定义。
在 Sarsa 中，策略通常是 $\varepsilon$-贪心的，这同样涉及最大化操作。
在这些算法中，会使用对最大值的估计作为最大值本身的估计，这可能导致显著的正偏差。
要明白原因，考虑一个状态 $s$，其中有许多动作 $a$，它们的真实价值 $q(s,a)$ 全都是零，但估计价值 $Q(s,a)$ 不确定，因此有的分布在零以上，有的分布在零以下。
真实价值的最大值是零，但估计值的最大值是正的，也就是存在正偏差。
我们把这称为\textbf{最大化偏差}。

\noindent\textbf{例 6.7：最大化偏差例子}
图 6.5 插入处所示的小 MDP，提供了一个简单例子，说明最大化偏差怎样伤害 TD 控制算法的表现。
这个 MDP 有两个非终止状态 $A$ 和 $B$。
一幕总是从 $A$ 开始，在 $A$ 中可以选择两个动作 left 和 right。
right 动作转移到终止状态，奖励和回报为零。
left 动作转移到 $B$，同样奖励为零；在 $B$ 中，有许多可能动作，所有动作都会立即终止，奖励从均值为 $-0.1$、方差为 1.0 的正态分布抽取。
因此，从 left 开始的任意轨迹期望回报都是 $-0.1$，所以在状态 $A$ 中选择 left 总是错误的。
尽管如此，我们的控制方法可能因为最大化偏差而偏好 left，使 $B$ 看起来具有正价值。
图 6.5 显示，带 $\varepsilon$-贪心动作选择的 Q-learning 一开始学到的 left 动作价值强到足以产生这种效果。
即使渐近地，Q-learning 选择 left 的次数仍比在我们参数设置下的最优次数多约 5\%（$\varepsilon=0.1$，$\alpha=0.1$，$\gamma=1$）。

\begin{figure}[htbp]
\centering
\includegraphics[width=0.74\textwidth]{tex/强化学习/images/math-7/math-7-fig-6-5.png}
\caption{在一个简单分幕 MDP（插入图）上比较 Q-learning 和 Double Q-learning。Q-learning 一开始学习到的 left 动作多于 right 动作，并且总是选择 left 的次数显著高于 $\varepsilon$-贪心动作选择所强制的 5\% 最小概率，其中 $\varepsilon=0.1$。相反，Double Q-learning 基本不受最大化偏差影响。这些数据是超过 10000 次运行的平均。初始动作价值全为零，$\varepsilon$-贪心动作选择中的平局随机打破。}
\end{figure}

是否存在能避免最大化偏差的算法？
作为起点，考虑一个赌博机情形：我们有许多对每个动作价值的带噪估计，它们是从每个动作的所有游玩所收到奖励的平均值。
如上所述，如果用估计值的最大值作为真实最大值的估计，就会产生正的最大化偏差。
一种看待这一问题的方式是，使用同一批样本（游玩）同时确定最大化动作并估计其价值。
假设我们把游玩随机分成两组，并用它们学习两个独立估计，称为 $Q_1(a)$ 和 $Q_2(a)$，每一个都是真实价值 $q(a)$ 的估计，对所有 $a\in\mathcal{A}$。
于是可以用一个估计，比如 $Q_1$，来确定最大化动作 $A^*=\arg\max_a Q_1(a)$，再用另一个估计 $Q_2$ 来提供其价值估计，即 $Q_2(A^*)=Q_2(\arg\max_a Q_1(a))$。
这个估计在无偏意义下是无偏的，即 $\mathbb{E}[Q_2(A^*)]=q(A^*)$。
也可以反过来重复这个过程，用两个估计的角色互换，得到第二个无偏估计 $Q_1(\arg\max_a Q_2(a))$。
这就是\textbf{双重学习}的思想。
注意，虽然我们学习两个估计，但每次游玩只更新其中一个估计；双重学习使内存需求加倍，但不会增加每一步的计算量。

双重学习思想可以自然扩展到完整 MDP 的算法。
例如，类似于 Q-learning 的双重学习算法称为 Double Q-learning，它把时间步分成两份，也许在每一步抛硬币决定。
如果硬币为正面，则更新为
\begin{equation}
Q_1(S_t,A_t)\leftarrow Q_1(S_t,A_t)+\alpha\left[
R_{t+1}+\gamma Q_2\!\left(S_{t+1},\arg\max_a Q_1(S_{t+1},a)\right)-Q_1(S_t,A_t)
\right].
\end{equation}
如果硬币为反面，则使用 $Q_1$ 与 $Q_2$ 互换的同样更新，用来更新 $Q_2$。
于是两个近似价值函数会被完全对称地处理。
行为策略可以同时使用两个动作价值估计。
例如，Double Q-learning 的 $\varepsilon$-贪心策略可以基于两个动作价值估计的平均值（或和）。
下面给出了 Double Q-learning 的完整算法。
在产生图 6.5 结果的例子中，双重学习似乎消除了最大化偏差造成的损害。
当然，也存在 Sarsa 和 Expected Sarsa 的双重版本。

\begin{center}
\fbox{\begin{minipage}{0.88\textwidth}
\noindent\textbf{Double Q-learning，用于估计 $Q_1\approx Q_2\approx q_*$}

\vspace{0.4em}
\noindent 算法参数：步长 $\alpha\in(0,1]$，小的 $\varepsilon>0$。

\noindent 初始化 $Q_1(s,a)$ 和 $Q_2(s,a)$，对所有 $s\in\mathcal{S}^+$、$a\in\mathcal{A}(s)$，且 $Q(terminal,\cdot)=0$。

\noindent 对每一幕循环：

\quad 初始化 $S$。

\quad 对一幕中的每一步循环：

\qquad 使用在 $Q_1+Q_2$ 上 $\varepsilon$-贪心的策略，从 $S$ 选择 $A$。

\qquad 采取动作 $A$，观察 $R,S'$。

\qquad 以 $0.5$ 的概率：
\[
Q_1(S,A)\leftarrow Q_1(S,A)+\alpha\left(R+\gamma Q_2(S',\arg\max_a Q_1(S',a))-Q_1(S,A)\right);
\]
\qquad 否则：
\[
Q_2(S,A)\leftarrow Q_2(S,A)+\alpha\left(R+\gamma Q_1(S',\arg\max_a Q_2(S',a))-Q_2(S,A)\right).
\]
\qquad $S\leftarrow S'$。

\quad 直到 $S$ 是终止状态。
\end{minipage}}
\end{center}

\noindent\textbf{练习 6.13}
对于带 $\varepsilon$-贪心目标策略的 Double Expected Sarsa，更新方程是什么？

\subsection*{6.8\quad 游戏、后状态和其他特殊情形}

在本书中，我们试图为一大类任务提供一种统一的方法；但当然，总有一些特殊任务最好用专门方式处理。
例如，我们的一般方法涉及学习一个\textbf{动作价值}函数，但在第 1 章中，我们给出了用于学习井字棋的 TD 方法，它学到的东西更像一个\textbf{状态价值}函数。
如果仔细看那个例子，就会发现所学函数既不是通常意义上的动作价值函数，也不是状态价值函数。
传统的状态价值函数评价智能体有选择动作余地的状态，但井字棋中使用的状态价值函数评价的是智能体已经走完一步之后的棋盘局面。
我们把这些称为\textbf{后状态}（afterstates），对应的价值函数称为\textbf{后状态价值函数}。
后状态在已经知道环境动态初始部分、但不一定知道完整动态的情形中很有用。
例如，在游戏中，我们通常知道自己走棋的即时效果。
我们知道每个可能的棋步会得到怎样的局面，但不知道对手会怎样回应。
后状态价值函数是一种自然方式，用来利用这类知识，从而产生更高效的学习方法。

用后状态来设计算法在井字棋例子中更高效，其原因很明显。
传统动作价值函数会把位置和走法的映射对应到价值估计。
但许多位置--走法对会产生相同的结果位置，如下例所示：

\begin{figure}[htbp]
\centering
\includegraphics[width=0.60\textwidth]{tex/强化学习/images/math-7/math-7-afterstates.png}
\end{figure}

在这类情形中，位置--走法对不同，但产生相同的“后位置”，因此必须具有相同价值。
传统动作价值函数必须分别评估这两个位置--走法对，而后状态价值函数会立即同时评估二者。
任何关于左侧位置--走法对的学习，都会立即转移到右侧那一对。

后状态不仅出现在游戏中，也出现在许多任务中。
例如，在排队任务中，存在把顾客分配给服务器、拒绝顾客或丢弃信息等动作。
在这类情形中，动作的直接效果事实上是完全已知的，而后续影响才完全未知。

不可能描述所有可能的专门问题类型及其对应的专门学习算法。
不过，本书发展出来的原则应当可以广泛适用。
例如，后状态方法仍然适当地描述为广义策略迭代，只是策略和（后）状态价值函数以本质上相同的方式相互作用。
在许多情况下，仍然要在同策略方法和离策略方法之间做选择，用来管理持续探索的需要。

\noindent\textbf{练习 6.14}
请说明 Jack 租车任务（例 4.2）怎样可以用后状态重新表述。
就这个具体任务而言，为什么这样的重新表述可能会加快收敛？

\subsection*{6.9\quad 总结}

在本章中，我们介绍了一种新的学习方法：时序差分（TD）学习，并展示了它怎样能够应用于强化学习问题。
和往常一样，我们把整体问题分为预测问题和控制问题，并且分别处理。
TD 方法是求解预测问题时蒙特卡洛方法的替代。
在两种情形中，扩展到控制问题都通过广义策略迭代（GPI）来实现，这是我们从动态规划中抽象出来的思想。
它的想法是，近似策略和价值函数应当以这样的方式相互作用：二者都朝各自的最优值移动。

构成 GPI 的两个过程之一，会驱动价值函数准确预测当前策略的回报；这是预测问题。
另一个过程会驱动策略相对于当前价值函数局部地改进（例如变成 $\varepsilon$-贪心）。
当第一个过程基于经验时，会出现一个复杂问题：如何维持足够探索。
我们可以按处理这个复杂问题的方式来分类 TD 控制方法：使用同策略方式，还是离策略方式。
Sarsa 是一种同策略方法，而 Q-learning 是一种离策略方法。
Expected Sarsa 也可以是离策略方法，正如我们在这里介绍的那样。
还有第三种方式可以把 TD 方法扩展到控制，本章没有包括，称为 actor--critic 方法。
这些方法将在第 13 章完整讨论。

本章介绍的方法，是如今使用最广泛的强化学习方法。
这大概是由于它们非常简单：它们可以在线应用，计算量很小，基于与环境交互产生的经验；它们几乎可以完全用单个方程表达，并能用很小的计算机程序实现。
在接下来的几章中，我们会扩展这些算法，使它们稍微更复杂，但也显著更强大。
所有新算法都将保留这里介绍的方法的本质：它们都能在线处理经验，计算量相对较小，并由 TD 误差驱动。
本章介绍的 TD 方法的特殊情形，应恰当地称为\textbf{一步、表格型、无模型} TD 方法。
在接下来的两章中，我们先把它们扩展到 $n$ 步形式（与蒙特卡洛方法相连），并进一步构造环境模型（与规划和动态规划相连）。
然后，在本书第二部分中，我们把它们扩展到各种形式的函数近似，而不是表格（与深度学习和人工神经网络相连）。

最后，本章一直在强化学习问题语境中讨论 TD 方法，但 TD 方法实际比这更一般。
它们是通用方法，用来学习对动态系统的长期预测。
例如，TD 方法可能与预测金融数据、寿命、选举结果、天气模式、动物行为、发电站负荷或顾客购买有关。
并且，TD 方法只与序列数据有关系；只要其理论性质首先得到充分理解，它们就可以独立于强化学习中的使用而广泛应用。
即便如此，TD 学习方法的这些其他潜在应用还没有得到广泛探索。

\section*{书目与历史注释}

正如我们在第 1 章中概述的，TD 学习思想的早期根源在动物学习心理学和人工智能中，最突出的包括 Samuel（1959）和 Klopf（1972）的工作。
Samuel 的工作在第 16.2 节中作为案例研究描述。
与 TD 学习相关的还有 Holland（1975，1976）关于价值预测之间一致性的早期思想。
这些思想影响了作者之一（Barto）；他在 1970 年到 1975 年是密歇根大学的研究生，而 Holland 当时在那里任教。
Holland 的思想导致了一些与 TD 有关的系统，包括 Booker（1982）的工作，以及 Holland（1986）的 bucket brigade；后者与下文讨论的 Sarsa 有关。

\noindent\textbf{6.1--2}
这些节中的大多数具体材料来自 Sutton（1988），包括 TD(0) 算法、随机游走例子，以及“时序差分学习”这个术语。
对其与动态规划和蒙特卡洛方法关系的刻画，受到 Watkins（1989）、Werbos（1987）等人的影响。
备份图的使用在本书第一版中是新的。

表格型 TD(0) 的均值收敛由 Sutton（1988）证明，概率 1 收敛由 Dayan（1992）证明，后者基于 Watkins and Dayan（1992）的工作。
这些结果由 Jaakkola, Jordan, and Singh（1994）以及 Tsitsiklis（1994）使用随机逼近已有强大理论的扩展加以扩充和加强。
其他扩展和推广将在后续章节中讨论。

\noindent\textbf{6.3}
Sutton（1988）确立了批量训练下 TD 算法的最优性。
Barnard（1993）把 TD 算法推导为两个步骤的组合：一步增量式地学习马尔可夫链模型，另一步根据该模型计算预测；这很好地说明了这个结果。
术语 certainty equivalence 来自自适应控制文献（例如 Goodwin and Sin, 1984）。

\noindent\textbf{6.4}
Sarsa 算法由 Rummery and Niranjan（1994）引入。
他们把它与人工神经网络结合考察，并称之为“Modified Connectionist Q-learning”。
“Sarsa”这个名称由 Sutton（1996）引入。
一步表格型 Sarsa（也就是本章处理的形式）的收敛性，由 Singh, Jaakkola, Littman, and Szepesvári（2000）证明。
“有风网格世界”例子由 Tom Kalt 建议。

Holland（1986）的 bucket brigade 思想，演化成一种与 Sarsa 密切相关的算法。
bucket brigade 的原始思想包含相互触发的规则链；它关注的是把信用从当前规则传回触发它的规则。
随着时间推移，bucket brigade 变得更像 TD 学习，会把信用传回任何时间上先前的规则，而不只是触发当前规则的那些规则。
Wilson（1994）详细说明了 bucket brigade 的现代形式在若干自然简化后几乎等同于一步 Sarsa。

\noindent\textbf{6.5}
Q-learning 由 Watkins（1989）引入；Watkins and Dayan（1992）把他给出的收敛证明大纲严格化。
更一般的收敛结果由 Jaakkola, Jordan, and Singh（1994）以及 Tsitsiklis（1994）证明。

\noindent\textbf{6.6}
Expected Sarsa 算法由 George John（1994）引入，他称之为“Q-learning”，并强调它作为离策略算法相对于 Q-learning 的优势。
当我们在本书第一版中把 Expected Sarsa 作为练习给出时，并不知道 John 的工作；van Seijen, van Hasselt, Whiteson, and Wiering（2009）在确立 Expected Sarsa 的收敛性质及其优于普通 Sarsa 和 Q-learning 的条件时，也不知道这项工作。
我们的图 6.3 改编自他们的结果。
Van Seijen 等人把“Expected Sarsa”专门定义为同策略方法（和我们在第一版中一样），而现在我们把这个名称用于目标策略和行为策略可以不同的一般算法。
Expected Sarsa 的一般离策略视角由 van Hasselt（2011）指出，他称之为“General Q-learning”。

\noindent\textbf{6.7}
最大化偏差和双重学习由 van Hasselt（2010，2011）引入并作了广泛研究。
图 6.5 中的示例 MDP 改编自他论文中的图 4.1（van Hasselt, 2011）。

\noindent\textbf{6.8}
后状态的概念与“决策后状态”（post-decision state）的概念相同（Van Roy, Bertsekas, Lee, and Tsitsiklis, 1997；Powell, 2011）。

\end{document}
